\chapter{Pruebas y Validación de Producto}
\label{ch:pruebas}

Este capítulo cubre cómo se testea \textbf{Seshat} y cómo se pone en manos de sus usuarios. La
comprobación se apoya en dos niveles complementarios: una suite de pruebas automatizadas
—unitarias y de integración— verifica que el código se comporta como se espera, y el \textit{gate}
de evaluación (capítulo~\ref{ch:evaluacion}) protege la calidad de los agentes frente a regresiones.
El tercer elemento es de otra naturaleza: una interfaz web que ofrece a los usuarios una forma
sencilla de operar el sistema —enviar reuniones, revisar los nodos extraídos y explorar el grafo—
sin tener que construir a mano las peticiones, extensas y anidadas, que la \ac{API} espera y necesita.
Este capítulo recorre los tres.

\section{Tests unitarios y de integración}
\label{sec:tests}

La suite de pruebas separa lo que puede verificarse de forma aislada y barata de lo que exige
levantar sistemas externos. Esa separación no es sólo organizativa: gobierna qué se ejecuta en cada
contexto —en cada \textit{commit}, en cada \textit{pull request} o sólo bajo demanda— y es la base
de la canalización de integración continua descrita en la sección~\ref{sec:ci}.

\subsection{Pruebas unitarias}

Las pruebas unitarias cubren la lógica pura y los adaptadores con todo el mundo exterior sustituido
por \textit{mocks}: el cálculo de confianza heurística, los validadores de entrada, la fusión de
configuración, los \textit{scorers} de evaluación, la lógica de los \textit{routers} de la \ac{API}
—que se prueban sin levantar un servidor, sustituyendo las dependencias inyectadas y ejercitando la
aplicación en proceso— o el comportamiento base de los agentes. Son rápidas y deterministas, no necesitan credenciales ni
servicios, y constituyen el grueso de la suite. Los \textit{mocks} se apoyan en las interfaces
abstractas que aíslan a \textbf{Seshat} de sus proveedores: los \ac{LLM} se sustituyen por
\textit{mocks} que devuelven salidas estructuradas predefinidas, y componentes como el
\textit{reranker} cuentan con implementaciones falsas de su interfaz para probar la lógica que los
rodea.

\subsection{Pruebas de integración}

Las pruebas de integración ejercitan el código contra las dependencias reales —PostgreSQL con
pgvector, el almacenamiento S3 y Secrets Manager emulados por LocalStack, y MLflow— levantadas
mediante los mismos contenedores de Docker Compose que se usan en desarrollo, sin réplicas ni
simuladores de esas dependencias. Cada prueba parte de un estado limpio: la base de datos de prueba se crea y
migra una vez, y el estado se reinicia entre pruebas vaciando las tablas y recreando la colección
vectorial. No todas requieren una clave de un proveedor de pago: la persistencia en pgvector, por
ejemplo, se valida con un modelo de \textit{embeddings} determinista que evita llamar a un servicio
real.

El eje que estructura estas pruebas es un sistema de \textit{marcadores} de \texttt{pytest}
organizados de forma jerárquica: \texttt{integration} engloba toda prueba que necesita un servicio
externo, y dentro de ella \texttt{llm} —y sus subconjuntos \texttt{agents}, \texttt{embedding},
\texttt{transcription} y \texttt{reranker}— señala las que requieren una clave de un proveedor en
vivo. Por defecto, la configuración de \texttt{pytest} excluye el marcador \texttt{llm}, de modo que
la ejecución habitual es rápida y gratuita; las pruebas que consumen llamadas de pago a los modelos
son explícitamente opcionales. Este diseño responde a una tensión real: las pruebas más valiosas
para un sistema basado en \ac{LLM} son también las más lentas, caras y no deterministas, y forzar su
ejecución en cada \textit{commit} sería inviable.

Las pruebas dependientes de un servicio están protegidas además por una doble compuerta: al
marcador se suma una condición de salto (\texttt{skipif}) que sondea la disponibilidad real del
servicio o de la clave. Así, un mismo fichero de prueba se comporta correctamente tanto en la
integración continua —con los servicios levantados— como en una máquina de desarrollo sin
credenciales, donde simplemente se omite en lugar de fallar. Para las pruebas que sí necesitan un
modelo, un ayudante selecciona automáticamente el proveedor disponible más económico, lo que hace
la suite portable a cualquier conjunto de credenciales del ejecutor.

La prueba de mayor alcance es un \textit{golden path} de extremo a extremo que recorre el
\textit{pipeline} completo —ingesta, transcripción, identificación y resolución— sobre una
transcripción breve, con Postgres, LocalStack y un modelo real: la aproximación más cercana a una
prueba de sistema.

\section{Eval gate y regresión automatizada}
\label{sec:eval-regresion}

Las pruebas de la sección anterior verifican el comportamiento del código, pero no bastan para un
sistema cuyo núcleo es un modelo de lenguaje: un cambio de \textit{prompt}, de modelo o de corpus
puede degradar la calidad de la extracción sin romper ninguna aserción. Esa clase de regresión es
justamente la que el \textit{gate} de evaluación vigila. Sus cinco arneses, su corpus y sus
umbrales se describen en el capítulo~\ref{ch:evaluacion}; lo que aquí importa es su papel como
mecanismo de control de regresión.

El \textit{gate} funciona como una prueba de aceptación de la calidad de los agentes, no del código.
Al ejecutarse contra el corpus de referencia produce un veredicto agregado —aprobado o suspenso
frente a los umbrales fijados— que se persiste en un fichero versionable. Comparar ese fichero entre
dos ejecuciones convierte la calidad en una magnitud observable: si una modificación hace caer la
precisión de un tipo de concepto por debajo de su umbral, el \textit{gate} pasa a suspenso y la
regresión se detecta antes de llegar a producción. Ese veredicto tiene consecuencias operativas: la
\ac{API} comprueba el \textit{gate} en el arranque y se niega a servir si no está en verde
(sección~\ref{sec:orquestacion}), de modo que la calidad medida es una precondición del despliegue,
no un informe que se pueda ignorar.

A diferencia de las pruebas unitarias y de integración, el \textit{gate} no forma parte de la
canalización de integración continua: reproducir sus llamadas reales a los modelos en cada
\textit{commit} sería caro y no determinista, por las mismas razones que apartan a las pruebas
\texttt{llm} de la ejecución por defecto. Se ejecuta, por tanto, de forma deliberada —antes de una
fusión sensible, de un cambio de modelo o de un despliegue— y su resultado se conserva junto al
código como constancia de la calidad con la que se validó cada versión.

\section{Experiencia de usuario (UX)}
\label{sec:ux}

La \ac{API} REST es el único punto de entrada al sistema, pero operarla directamente exige construir
peticiones HTTP a mano —subir un archivo por \textit{multipart}, componer el cuerpo JSON de una
revisión con las decisiones de cada nodo, encadenar consultas para recorrer el grafo—. Para que un
usuario pueda trabajar con \textbf{Seshat} sin ese esfuerzo, el proyecto incluye una interfaz web:
una herramienta de conveniencia, una capa de ergonomía sobre la \ac{API} que traduce esas operaciones en
formularios, botones y vistas, y hace el sistema utilizable en la práctica sin conocer la \ac{API}.

Su rasgo de diseño más relevante es precisamente lo que \emph{no} hace. La interfaz se comunica con
el sistema exclusivamente a través de la \ac{API}: no accede a la base de datos, al almacenamiento de objetos
ni a ningún servicio interno, y no reimplementa ninguna regla del \textit{pipeline}. La aprobación
automática, la puntuación de confianza, la resolución de relaciones y las transiciones de estado ocurren
todas en el servidor; la interfaz sólo dispara acciones y representa lo que la \ac{API} le devuelve. Es,
además, sin estado: se autentica introduciendo una clave de \ac{API} —la misma que gobierna los roles
(sección~\ref{sec:orquestacion})— y no persiste nada entre sesiones. Construida sobre Streamlit,
esta delegación estricta la mantiene como un cliente fino que no puede divergir del comportamiento
del sistema.

Sobre esa base, la interfaz cubre el ciclo de vida completo de un \textit{pipeline}. Un operador
envía un archivo desde un formulario y sigue su avance por la máquina de estados en una barra de
progreso que sondea periódicamente el estado del trabajo; no gobierna las transiciones —de eso se
encarga el \textit{worker}— sino que las refleja. Cuando el trabajo se detiene en revisión, presenta
la interfaz de \acs{HITL} descrita en el capítulo~\ref{ch:implementacion}: aprobación masiva por
umbral de confianza y decisiones individuales de aprobar, rechazar o editar cada nodo, con acceso al
fragmento de transcripción que respalda cada uno. Fuera del flujo de revisión, ofrece acciones
manuales —añadir nodos o relaciones, corregir contenido mediante anulaciones justificadas— y
herramientas de exploración del grafo: consulta por filtros, búsqueda semántica y una vista
interactiva del impacto de un nodo sobre sus vecinos. Un panel de administración, protegido por la
clave raíz, gestiona las claves de \ac{API}. En conjunto, la interfaz hace visible y operable desde el
navegador lo que la \ac{API} expone de forma programática, sin alterar lo que ya se ha descrito en los
capítulos anteriores.
